\section{Related work} 
\label{sec:rw}
\textbf{Stochastic Adaptive First-order Optimizers in Deep Learning.} The idea of adaptive first-order optimizers dates back to RProp~\citep{braun1992rprop}. AdaGrad~\citep{duchi2011adaptive} adapted the learning rate of features by estimated geometry and assign larger learning rate to infrequent features. 
RMSProp~\citep{hinton2012neural} generalized RProp and is capable to work with smaller batch sizes. Adam~\citep{kingma2014Adam} improved RMSProp by introducing a running average of gradients, and has so far become the dominant approach to solve optimization problems in deep learning, especially for training Transformers~\citep{vaswani2017attention}. Many follow-up works proposed variants of Adam~\citep{dozat2016incorporating,shazeer2018adafactor,reddi2019convergence,loshchilov2017decoupled, zhuang2020adabelief, you2019large}. \citet{chen2023symbolic} performed a search over adaptive first-order algorithms and discovered Lion, which is a improved version of sign momentum SGD.


\textbf{Second-order Optimizers in Deep Learning.} Second-order optimizers are believed to have the potential to outperform adaptive first-order optimizers. Classical second-order optimization algorithms pre-condition the gradient with curvature information~\citep{broyden1970convergence,nesterov2006cubic,conn2000trust}. Over the years, people have developed numerous ways to adapt these methods to deep learning. To the best of our knowledge, \citet{becker1988improving} was the first to use diagonal Hessian as the pre-conditioner. \citet{martens2010deep} approximated the Hessian with conjugate gradient. \citet{schaul2013no} automatically tuned learning rate of SGD by considering diagonal Hessian. \citet{pascanu2013revisiting} considered Gaussian Newton's approximation of Hessian and Fisher information matrix. \citet{martens2015optimizing} and follow-up works \citep{ba2017distributed,george2018fast,martens2018kronecker,zhang2022eva} proposed to approximate the Hessian based on the structure of neural networks. \citet{yao2021adahessian, jahani2021doubly} proposed to use the EMA of diagonal Hessian estimator as the pre-conditioner. 

Despite these progress on deep learning applications, for decoder-only large language models, Adam still appears to the most popular optimizer. The authors of this paper suspect that many previous second-order optimizers face the challenge that the computational / memory overhead due to frequent Hessian computation hinders improvements in wall-clock time~\citep{martens2015optimizing,gupta2018shampoo}. Some of them also depend on specific model architecture or hardware structures, e.g., \citet{anil2020scalable} offloads hessian computation to CPUs, and \citet{george2018fast} needs ResNets and very large batch size to approximate the Fisher information matrix. To the best of our knowledge, there was no previous report that second-order optimizers can achieve a speed-up on large language models in total compute.

\textbf{Gradient Clipping.} Global gradient clipping has been a standard practice in pre-training language models~\citep{merity2017regularizing,radford2019language,izsak2021train,zhang2022opt}. It helps stabilizes training and avoids the effect of rare examples and large gradient noise. \citet{zhang2019gradient,mai2021stability} showed that global gradient clipping is faster than standard SGD when global smoothness does not hold. \citet{zhang2020adaptive,crawshaw2022robustness} found out per-coordinate gradient clipping can function as adaptivity. In addition to gradient clipping, \ours~is the first to clip the update (coordinate-wise) in second-order methods to avoid the effect of Hessian's changing along the trajectory and the inaccuracy of Hessian approximation. 

\textbf{Optimization Algorithms in LM Pre-training.} Adam~\citep{kingma2014Adam} (with decoupled weight decay~\citep{loshchilov2017decoupled}) has become the dominant approach for language model pre-training~\citep{vaswani2017attention,devlin2018bert,radford2019language,brown2020language,zhang2022opt,touvron2023llama}. Different from vision tasks with CNNs~\citep{he2016deep} where models trained with SGD generalize better than models trained with Adam, Adam outperforms SGD by a huge margin on language modeling tasks with Transformers~\citep{anil2019memory,liu2020understanding,kunstner2023noise}. \citet{raffel2020exploring,chowdhery2022palm} trained Transformers with AdaFactor~\citep{shazeer2018adafactor}, which is a low rank version of Adam. \citet{you2019large} proposed to make the update of Adam proportional to per-layer paramter norm to stably train LLMs. 

\section{Conclusion}
We introduced Sophia, a scalable second-order optimizer for language model pre-training. Sophia converges in fewer steps than first-order adaptive methods, while maintaining almost the same per-step cost. On language modeling with GPT models, Sophia achieves a 2x speed-up compared with AdamW in the number of steps, total compute, and wall-clock time.

\subsection*{Acknowledgements}

We thank Jeff Z. HaoChen, Neil Band, Garrett Thomas for valuable feedbacks.
HL is supported by Stanford Graduate Fellowship. The authors
would like to thank the support from NSF IIS 2211780.